\begin{abstract}
This study asks how much of a retrieval gain reaches the final answer of a retrieval-augmented generation (RAG) system.
We are not aware of work that measures retrieval-to-answer transfer ratios jointly with measured cost on free hardware.
We fix the reader (Qwen2.5-7B-Instruct-AWQ) and the corpus (a 300,000-passage BEIR-NQ sub-corpus) and vary only the retriever: BM25 and six open-source encoders from 22.7M to 595.8M parameters.
We measure encoding time, query latency, index size and GPU memory on one free T4.
On 2,255 Natural Questions queries, Hit@5 ranges from 66.4\% to 81.6\% across the dense encoders, but exact match (EM) ranges only from 47.1\% to 51.0\%.
Relative to all-MiniLM-L6-v2, the transfer ratio $\Delta\mathrm{EM}/\Delta\mathrm{Hit@5}$ lies between 0.19 and 0.25 for the other five dense encoders.
$\sim$20--25\% of Hit@5 gain appears as EM gain, within our setup.
bge-large-en-v1.5 and Qwen3-Embedding-0.6B need 3--7.5$\times$ the T4 encoding time of e5-base-v2.
Both retrieve worse than e5-base-v2 on Hit@5 (exact McNemar test) and nDCG@10 (Wilcoxon signed-rank test), significant after Holm correction over 21 pairs.
Both show no detectable EM gain over e5-base-v2.
\end{abstract}
